\documentclass{article}
\usepackage[margin=1in]{geometry}
\usepackage{array}
\usepackage{amsmath}
\usepackage{fancyvrb}
\usepackage{caption}
\captionsetup[table]{skip=10pt}

\begin{document}
\section{Confidence Intervals, etc.}
If a value $z_c$ or $t_c$ value with $N-1$ degrees of freedom corresponds to our desired confidence interval:
\begin{equation*}
    \mu \in \left(\bar{X} -\frac{z_c\sigma}{\sqrt{N}},\bar{X}+\frac{z_c\sigma}{\sqrt{N}}\right)
\end{equation*} 
If it is a one tailed test, we use $\alpha=1-(confidence_{interval})$, but for a two-tailed test, we use $\alpha = \frac{\alpha_{onetail}}{2}$
\section{RegEx}
In other document.
\section{Linear Regression}

\begin{equation*}
    \beta = Vy
\end{equation*}
Where 
\begin{equation*}
    V = (X^T X)^{-1}X^T
\end{equation*}
Assume $\beta = [a_1,a_2,\ldots, b]^T$. \\
\begin{equation*}
    \hat{y} = a_1 x_1 +a_2 x_2 + \cdots + b
\end{equation*}
If there are N samples, with M features + bias term (b).\\
Without bias term:
\begin{eqnarray*}
    \dim(X) = N\times M\\
    \dim(y) = N\times 1\\
    \dim(\beta) = M\times 1 
\end{eqnarray*}
That is X has N rows, M columns.\\
With bias term:
\begin{eqnarray*}
    \dim(X) = N\times (M+1)\\
    \dim(y) = N\times 1\\
    \dim(\beta) = (M+1)\times 1 
\end{eqnarray*}
Essentially, the bias term is an additional feature when looking at dimensions.\\
If there is no bias term initially, or provided, one can assume all bias terms have a value of 1, making it an additional last column to the X matrix.
\subsection{Normalization}
We calculate $\mu$ and $\sigma$ for each column in $X$ and $y$ where 
\begin{eqnarray*}
    \mu_n = \frac{\sum x}{\text{len($x_n$)}}\\
    \sigma_n = \sqrt{\frac{\sum (x-\mu)^2}{\text{len($x_n$)}}}
\end{eqnarray*}
Normalize using:
\begin{equation*}
    z=\frac{x-\mu}{\sigma}
\end{equation*}
\subsection{MSE}
\begin{equation*}
    MSE = \frac{\sum E^2}{N}
\end{equation*}
where $E$ is the error of a sample, and $N$ is the number of samples.
\section{Ridge Regression}
\begin{equation*}
    \beta = Vy
\end{equation*}
Where 
\begin{equation*}
    V = (X^T X+\lambda I)^{-1}X^T
\end{equation*}
If we want to perform N-fold cross-validation, and we are considering M values of $\lambda$, there are $MN$ possibilities of trainable models.\\
The larger the value of $\lambda$, the less it fits the model. So, a very small lambda is likely to overfit, or, at the very least, allow for overfitting.\\
Ridge Regression Loss is:
\begin{equation*}
    \sum(y-\hat{y})^2+\lambda\sum \beta^2
\end{equation*}
\section{N-gram}
\begin{itemize}
    \item word-based: N consecutive words, starting from one word after the previous one. Will have exactly N words.
    \item character-based: N consecutive characters, rest is basically the same as the above.
\end{itemize}
Term frequency:
\begin{equation*}
    TF(t,d) = \frac{\# t \text{ in } d}{\text{len($d$)}}
\end{equation*}
Inverse Document Frequency:
\begin{equation*}
    IDF(t,D) = \log\left(\frac{\#\text{documents}}{\#\text{documents with }t}\right)
\end{equation*}
Importance of word:
\begin{equation*}
    TF\text{-}IDF(t,d,D) = TF(t,d) \times IDF(t,D)
\end{equation*}
\end{document}
